Os ensaios foram planejados para avaliar a camada WCAN isoladamente, sem tornar
o barramento CAN físico o fator dominante da análise. A integração com TWAI foi
validada previamente em bancada com dois transceptores VP230, mas a campanha
principal concentrou-se no enlace sem fio, nas filas internas, na filtragem e nos
mecanismos de confirmação.

\section{Ambiente Experimental}

O ambiente experimental foi organizado em torno de um servidor Ubuntu usado como
ponto central de acesso remoto, orquestração e coleta dos testes. Essa escolha
facilitou o controle das placas espalhadas na bancada, a recuperação dos logs e a
execução dos scripts de automação sem depender de intervenção manual contínua.

A campanha utilizou cinco placas ESP32 conectadas por USB serial ao computador de
monitoramento, sendo três ESP32 e duas ESP32-C3. Em cada repetição, as placas
foram sorteadas como sensores, receptores ou nós ociosos, reduzindo a chance de
associar um resultado a uma placa específica. Cada ensaio durou 30 s, com três
repetições por topologia e 5 s de intervalo entre execuções.

O objetivo desse arranjo foi tornar a execução e a análise dos testes o mais
automáticas possível. Em vez de tratar cada ensaio como uma operação manual,
o ambiente foi estruturado para que a configuração, a sincronização dos papéis,
o registro das saídas e a posterior consolidação dos dados fossem conduzidos por
scripts.

\section{Automação dos Testes}

Os testes foram automatizados pelos scripts em \texttt{transmitter/testing}. O
arquivo \texttt{boards.yaml} descreve as placas disponíveis e o arquivo
\texttt{tests.yaml} define suítes, frequências, topologias, duração, repetições,
transporte, valor de \texttt{linger\_ms} e quantidade de CAN IDs por sensor. Antes
da execução, o runner gera \texttt{run\_plan.json} e, para cada teste, um
\texttt{scenario.json} com a topologia, os filtros e a linha de configuração
enviada a cada placa.

A \texttt{main.cpp} foi mantida como uma camada fina de inicialização e exemplo de
uso. Em vez de concentrar nela a lógica da automação de testes, o \texttt{app\_main} apenas
prepara o sistema, lê a configuração de boot e despacha a execução para módulos
específicos conforme o papel informado. A parametrização do ensaio fica isolada em
\texttt{runtime\_config.cpp}, enquanto o comportamento de sensor e receptor é
encapsulado, respectivamente, em \texttt{sensor\_app.hpp} e
\texttt{receiver\_app.hpp}. Essa separação permite reutilizar a mesma imagem de
firmware em diferentes papéis sem modificar o fluxo central, simplificando a
automação e a repetição dos testes.

Após gravar os firmwares necessários, o monitor serial reinicia as placas, aguarda
a mensagem \texttt{WCAN\_CFG\_WAIT} e envia por UART uma linha de configuração
específica para cada placa. Essa linha existe para que o firmware não precise ser
recompilado a cada repetição: ela informa o papel da placa, o transporte esperado, 
o CAN ID base do sensor, a quantidade de CAN IDs por sensor, a frequência de geração,
o valor de \texttt{linger\_ms} e, no caso de receptores, a lista de filtros de CAN ID. 
Em termos práticos, a mesma imagem de firmware pode atuar como sensor, receptor ou 
nó ocioso, e a distinção é feita apenas pelo conteúdo enviado nessa fase de boot. 
As saídas seriais são salvas separadamente por placa, o que permite reconstruir a 
sequência de eventos e inspecionar manualmente os casos reprovados.

\section{Matriz de Ensaios e Métricas}

A matriz principal de topologias incluiu todas as combinações com pelo menos um sensor e um 
receptor, limitadas pelo total de cinco placas disponíveis, isto é, \(S \geq 1\), \(R \geq 1\) 
e \(S + R \leq 5\). Isso resultou em dez topologias: 1S-1R, 1S-2R, 1S-3R, 1S-4R, 2S-1R, 2S-2R, 
2S-3R, 3S-1R, 3S-2R e 4S-1R.

A execução \texttt{full\_run} foi organizada em cinco suítes. A suíte \texttt{baseline} usou 
um CAN ID por sensor e \texttt{linger\_ms = 100}, servindo como referência de desempenho. 
A suíte \texttt{multiple} manteve \texttt{linger\_ms = 100}, mas configurou quatro CAN IDs por 
sensor, simulando um módulo de aquisição que publica vários sinais simultaneamente; nessa suíte 
também foi avaliada a variante multicast, adotada para múltiplos receptores interessados. A suíte 
\texttt{real\_time} configurou \texttt{linger\_ms = 0}, forçando cada amostra a ser enviada em 
seu próprio pacote ESP-NOW para medir o custo de operar sem batching. A suíte \texttt{active\_filter} 
usou pelo menos dois sensores e configurou allowlists nos receptores, validando a filtragem por 
CAN ID. Por fim, a suíte \texttt{mixed\_frequency} atribuiu frequências diferentes aos sensores 
de uma mesma topologia, sorteadas entre 100, 200, 750 e 1000 Hz, para verificar se fluxos rápidos 
prejudicam fluxos lentos.

\begin{table}[htbp]
\centering
\scriptsize
\setlength{\tabcolsep}{3pt}
\renewcommand{\arraystretch}{1.15}
\begin{tabular}{p{0.19\linewidth}|p{0.26\linewidth}|p{0.25\linewidth}|p{0.22\linewidth}}
\hline
Suíte & O que foi testado & Resultado esperado & Interpretação pretendida \\
\hline
\texttt{baseline} & Um CAN ID por sensor, \texttt{linger\_ms = 100}, varredura de 10 a 1000 Hz. & Entrega próxima de 100\% até frequências altas, pois o batching reduz a quantidade de pacotes. & Estabelecer a referência de desempenho do WCAN com batching. \\
\hline
\texttt{multiple} & Quatro CAN IDs por sensor, cada um com sua fila e sua task de processamento. & Aumentar a pressão sobre filas, timers e tarefas, mas sem colapsar se a arquitetura por CAN ID estiver correta. & Verificar concorrência entre fluxos do mesmo módulo. \\
\hline
\texttt{real\_time} & \texttt{linger\_ms = 0}, isto é, envio amostra a amostra. & Degradação em frequências altas, por aumento de overhead MAC e ocupação de filas. & Confirmar se batching é essencial para a proposta. \\
\hline
\texttt{active\_filter} & Receptores aceitando apenas subconjuntos dos CAN IDs transmitidos. & Receptores não devem entregar nem confirmar mensagens fora da allowlist. & Validar filtragem e comportamento de ACK com receptores parciais. \\
\hline
\texttt{mixed\_frequency} & Sensores na mesma topologia transmitindo em frequências diferentes. & Fluxos rápidos não devem impedir a entrega dos fluxos lentos. & Avaliar convivência entre sinais de naturezas diferentes. \\
\hline
\end{tabular}
\caption{Suítes executadas na campanha \texttt{full\_run} e objetivo de cada uma.}
\label{tab:suites-executadas}
\end{table}

\section{Métricas Coletadas}

A taxa de entrega foi calculada a partir dos valores de contador transportados pelos lotes. 
No lado sensor, o firmware registra uma linha \texttt{CAN\_PROC\_*} para cada lote colocado 
no transporte, contendo o CAN ID e o intervalo de contadores do lote, por exemplo \texttt{[0..59]}. 
No lado receptor, o callback \texttt{wcan\_recv\_callback} registra o mesmo intervalo quando 
o pacote é recebido. O script \texttt{analysis.py} reconstrói os intervalos recebidos e calcula 
a fração de contadores transmitidos que aparecem em pelo menos um receptor esperado. Na suíte 
de filtragem ativa, receptores cuja allowlist não contém determinado CAN ID não entram no 
denominador daquele fluxo.

Também foi analisada a entrega par a par entre cada sensor e cada receptor. Essa segunda métrica 
é informativa, mas não é o critério principal de sucesso do broadcast, porque o mecanismo de 
ACK dessa variante segue a semântica do CAN: se qualquer receptor interessado confirma o pacote, 
o transmissor considera a transmissão bem-sucedida. Portanto, um pacote pode ser aprovado na 
métrica principal e ainda assim não ter chegado a todos os receptores da topologia. Essa diferença 
é relevante para interpretar os resultados de topologias com múltiplos receptores.

A latência foi medida de formas diferentes nas estratégias de transporte, portanto os valores 
não devem ser comparados diretamente. No broadcast, a métrica \texttt{LAT\_RTT} registra o tempo 
de ida e volta, em milissegundos, até o recebimento do ACK de aplicação. No multicast, a métrica 
\texttt{LAT\_CB} registra, em microssegundos, o tempo entre a chamada a \texttt{esp\_now\_send()} 
e a execução do callback de transmissão do driver ESP-NOW. Esse segundo valor mede o retorno do driver 
e a confirmação no nível MAC, não a entrega fim a fim até o processamento da aplicação receptora.

O airtime foi estimado pela instrumentação do firmware. A cada janela periódica, o código acumula 
a quantidade de bytes transmitidos por ESP-NOW e, assumindo taxa física de 6 Mbps, estima a fração 
de tempo de canal consumida. O resultado é registrado nas linhas \texttt{MEASURE}, no campo 
\texttt{util\_per\_mille}, em partes por mil; 10 partes por mil correspondem a 1\% do tempo de canal. 
A estimativa ignora parte do overhead físico e deve ser interpretada como métrica relativa entre 
variantes.

Também foi coletado o high-water mark de pilha das tasks FreeRTOS. A instrumentação registra, nas 
linhas \texttt{TASK}, o campo \texttt{hwm\_bytes}, que indica o menor número de bytes livres observado 
na pilha daquela task. Valores abaixo de 100 bytes foram tratados como sinal de risco, pois indicam 
operação próxima de overflow de pilha. O analisador também verificou crashes, quedas de fila do 
produtor, quedas na fila interna do WCAN, falhas por limite de retry, lacunas em \texttt{RX\_STATS}, 
integridade de heap e validações específicas de cada suíte.
